% ============================================================================
% THE VIRTUAL MACHINE: TOKENIZE → STORE → FOLD → PROMPT → EXECUTE → REIFY
% ============================================================================
\subsection{Machine Architecture}
\label{sec:vm_machine}

The Machine is the top-level orchestrator of the entire architecture.  It
manages the lifecycle of a single data-to-inference pipeline by composing
four RPC-mediated subsystems into a unified processing loop:

\begin{enumerate}[nosep]
  \item \textbf{Loading:} Raw byte streams are tokenized via the Universal
        Tokenizer into Morton-packed keys.  Each key carries a position
        coordinate and a symbol byte, interleaved via a Z-order curve so that
        locality is preserved across scales.
  \item \textbf{Storage:} Keys are inserted into the DMT Forest (a radix
        trie) through \texttt{WriteBatch} RPCs, providing prefix-compressed
        persistent storage for the full tokenized corpus.
  \item \textbf{Graph Construction:} The same key stream is simultaneously
        routed to the Graph substrate, which converts it into
        \texttt{primitive.Value} signals and constructs a fold hierarchy via
        RecursiveFold (\Cref{sec:self_programming}).
  \item \textbf{Prompt Resolution:} Prompt strings are tokenized through the
        same pipeline, compiled into Value sequences, and dispatched to
        CantileverServer for multi-stage BVP resolution
        (\Cref{sec:cantilever_bvp}).
\end{enumerate}

All subsystem capabilities (Tokenizer, Forest, Graph, MacroIndex,
CantileverServer) are resolved at runtime through a cluster Router, which
performs load-balanced capability selection across registered servers.
Cap'n~Proto RPC provides the transport layer; every inter-service call is a
zero-copy capability invocation over in-process pipes.

\subsection{The Self-Programming Loop}
\label{sec:self_programming}

The central design contribution is a closed feedback loop in which the
system discovers, verifies, and internalizes its own computational operators.
The loop has five stages:

\subsubsection{Observe}

Data bytes are tokenized into Morton keys, stored in the radix trie, and
assembled into \texttt{primitive.Value} signals.  The Graph substrate
constructs a topology-guided fold hierarchy using RecursiveFold: all pairwise
Jaccard similarities are computed, the most similar pairs merge first, and the
shared structural invariant (bitwise AND) becomes a fold label while
directional residues (Hole) recurse to deeper levels.  The result is a
persistent reasoning graph where every node records a shared invariant and
two residue children.

\subsubsection{Prompt}

CantileverServer resolves prompts in two stages:

\begin{enumerate}[nosep]
  \item \textbf{Exact lexical continuation.}  The prompt is tokenized through
        the same Universal Tokenizer path, producing a Value sequence.  This
        sequence is decoded to bytes and matched against the stored corpus
        rows; the first exact prefix match yields the suffix as a
        continuation.
  \item \textbf{Operator-mediated bridging.}  When no exact prefix match
        exists, the system attempts BVP bridging.  For each corpus row longer
        than the prompt, the prompt's accumulated OR-fold signal is treated as
        the start boundary and the row's suffix signal as the goal boundary.
        The geometric \texttt{AffineKey} delta between start and goal is
        computed (one 64-bit word per Value block), and the MacroIndex is
        queried for an operator that resolves the gap.  The MacroIndex
        supports both exact lookup (by full AffineKey) and approximate lookup
        (via PhaseDial embeddings with cosine-similarity nearest-neighbor
        search among hardened opcodes).
\end{enumerate}

\subsubsection{Harden}

Successful Stage~2 bridges feed back into the MacroIndex via
\texttt{RecordCandidateResult}.  Each bridge outcome records pre-residue,
post-residue, an ``advanced'' flag (post-residue $<$ pre-residue), and a
``stable'' flag.  Candidates that succeed repeatedly---advanced, stable, and
residue-reducing for $\ge 5$ consecutive uses---are promoted from
\emph{Candidate} to \emph{Hardened} status.  Hardening triggers PhaseDial
embedding of the operator's AffineKey, making it available for approximate
lookup by future prompts that encounter geometrically similar gaps.

\subsubsection{Execute}

Hardened MacroOpcodes are first-class Values.  \texttt{EncodeMacroOperator}
produces a Value carrying \texttt{OpcodeMacro} in its opcode block, the
learned affine operator $f(x) = \mathit{scale} \cdot x + \mathit{translate}
\pmod{8191}$ in its shell, and an XOR-signature of the operator's key in its
core blocks.  These Values can be placed directly into program sequences and
executed by the Interpreter at native speed, without re-deriving the
transformation at runtime.

\subsubsection{Reify}

\texttt{ReifyTrace} closes the loop.  After the Interpreter executes a
program, the execution trace---a sequence of \texttt{ExecutionStep} records,
each carrying the program counter, opcode, phase-before, and
phase-after---is composed into a single affine macro operator.  Composition
chains each step's affine transform:
\begin{equation}
  (s_{\mathrm{comp}},\; t_{\mathrm{comp}})
  = \prod_{i=1}^{n} (s_i,\; t_i)
  \quad\text{where}\quad
  s' = s_{\mathrm{comp}} \cdot s_i \!\!\pmod{8191},\;
  t' = s_{\mathrm{comp}} \cdot t_i + t_{\mathrm{comp}} \!\!\pmod{8191}
  \label{eq:reify}
\end{equation}
The resulting Value is a candidate for insertion into the MacroIndex, where
it enters the hardening pipeline.  A successful trace between two boundary
states thus crystallizes into a single instruction that future programs can
invoke directly.

\subsection{The Interpreter}
\label{sec:interpreter}

The Interpreter is a register machine that executes programs encoded as
sequences of \texttt{primitive.Value}.  Each Value is simultaneously
instruction and operand: the opcode block (C5) carries the control-flow
instruction and branch metadata, the shell (C7) carries the device address
and affine operator, and the core blocks (C0--C4) carry the
$\operatorname{GF}(8191)$ state.

Execution maintains a single running phase $\phi \in \operatorname{GF}(8191)$
initialized to the multiplicative identity.  At each program counter step,
the current node's affine operator transforms the phase:
\begin{equation}
  \phi \;\leftarrow\; a \cdot \phi + b \pmod{8191}
\end{equation}
The opcode then determines the program counter advance:

\begin{center}
\begin{tabular}{ll}
  \toprule
  \textbf{Opcode} & \textbf{Semantics} \\
  \midrule
  \texttt{Next}    & $\mathit{pc} \leftarrow \mathit{pc} + 1$ \\
  \texttt{Jump}    & $\mathit{pc} \leftarrow \mathit{pc} + \Delta$ (relative offset stored in node) \\
  \texttt{Branch}  & Evaluate $N$ candidate branches; pick lowest-residue winner; $\mathit{pc}$++ \\
  \texttt{Reset}   & $\phi \leftarrow 1$;\; $\mathit{pc}$++ \\
  \texttt{Halt}    & Emit accumulated state; stop \\
  \texttt{Macro}   & Apply learned affine: $\phi \leftarrow s \cdot \phi + t \pmod{8191}$;\; $\mathit{pc}$++ \\
  \bottomrule
\end{tabular}
\end{center}

\texttt{OpcodeMacro} is the critical instruction for self-programmability: it
executes a hardened learned operator at native interpreter speed.  The macro's
scale and translate fields are extracted from the Value's shell via
\texttt{MacroAffine} and applied as a single Mersenne-reduced multiply-add.
Guard radii provide an optional early-termination mechanism: if the phase
transition magnitude exceeds the node's guard radius, execution halts and the
current state is emitted.

\subsection{Batch Transport}
\label{sec:batch_transport}

Data flows between Machine and its subsystems through \texttt{WriteBatch}
RPCs.  The Tokenizer, Forest, and Graph each expose a streaming
\texttt{WriteBatch} method that accepts bulk payloads in a single Cap'n~Proto
call, amortizing per-key RPC overhead.  The Machine first sends the full byte
stream to the Tokenizer via \texttt{WriteBatch}, drains the resulting Morton
keys, then fans the key stream out to both Forest and Graph through their
respective \texttt{WriteBatch} endpoints.  This bulk pipeline replaces
per-element streaming and ensures that ingestion throughput scales with
payload size rather than element count.
